% The buffer as a strip of slots, in the ordinary case and in the wrapped case.
% Cell i occupies x from i*8mm to (i+1)*8mm; the strip is 128 mm wide.
\begin{tikzpicture}[font=\sffamily\small]

% ================================================ row 1: the ordinary case
\node[font=\sffamily\bfseries\small, anchor=west] at (0mm,64mm)
  {Ordinary case: the occupied span does not cross the end of the array};
\foreach \i in {0,1,2,11,12,13,14,15}
  {\node[mmres, minimum width=8mm, text width=6mm, minimum height=7mm]
     at ({\i*8mm},43mm) {\tiny \i};}
\foreach \i in {3,...,10}
  {\node[mmram, minimum width=8mm, text width=6mm, minimum height=7mm]
     at ({\i*8mm},43mm) {\tiny \i};}

\node[font=\sffamily\scriptsize, align=center] at (92mm,57mm)
  {head \& MASK = 11\\{\tiny the next slot to be written}};
\draw[flow] (92mm,54.0mm) -- (92mm,50.5mm);
\draw[flow] (28mm,39.0mm) -- (28mm,42.5mm);
\node[font=\sffamily\scriptsize, align=center] at (28mm,36mm)
  {tail \& MASK = 3\\{\tiny the next slot to be read}};

\node[font=\sffamily\scriptsize, anchor=west, text width=130mm] at (0mm,28mm)
  {head = 0x0000011B, tail = 0x00000113. Neither is ever wrapped. used = head
   minus tail = 8, and the buffer is full when that difference equals 16, so no
   slot is sacrificed to tell full from empty.};

% ================================================ row 2: the wrapped case
\node[font=\sffamily\bfseries\small, anchor=west] at (0mm,20mm)
  {Wrapped case: the same arithmetic, and not a special case in the code};
\foreach \i in {3,...,10}
  {\node[mmres, minimum width=8mm, text width=6mm, minimum height=7mm]
     at ({\i*8mm},0mm) {\tiny \i};}
\foreach \i in {0,1,2,11,12,13,14,15}
  {\node[mmram, minimum width=8mm, text width=6mm, minimum height=7mm]
     at ({\i*8mm},0mm) {\tiny \i};}

\node[font=\sffamily\scriptsize, align=center] at (28mm,13mm)
  {head \& MASK = 3\\{\tiny still the next slot to be written}};
\draw[flow] (28mm,10.0mm) -- (28mm,7.5mm);
\draw[flow] (92mm,-4.0mm) -- (92mm,-0.5mm);
\node[font=\sffamily\scriptsize, align=center] at (92mm,-7mm)
  {tail \& MASK = 11\\{\tiny still the next slot to be read}};

\node[font=\sffamily\scriptsize, anchor=west, text width=130mm] at (0mm,-16mm)
  {head = 0x00000123, tail = 0x0000011B. used = 8 again. The mask moved the
   subscript across the end of the array and the code did not notice, which is
   the whole reason the capacity is a power of two.};
\end{tikzpicture}
